\section{Dual Variables for the Outer Loop: A Mapping from Cluster Resource Allocation}
\label{app:cvx_mapping}

The recipe at the core of CHEETAH, relax a hard discrete resource-allocation
problem to a convex one, read the dual variables, and use them to
guide the discrete choices, is not unique to accelerator co-design. The same
structure appears in large-scale cluster task placement \citep{cvxcluster2026},
which solves granular server allocation $100$--$1000\times$ faster by (i) relaxing
the binary task-to-server assignment and \emph{aggregating servers into shapes} so
the program scales with the number of shapes rather than servers, and (ii)
extracting the LP's dual variables as per-resource \emph{dual variables} and placing
tasks greedily by \emph{net utility} $=$ priority $-$ $\sum_{r}(\lambda_r
\times \text{demand}_r)$, re-evaluating as the load shifts. We already borrow the shape
aggregation (\cref{app:shape_agg}); this appendix records the full one-to-one
correspondence and, more usefully, the concrete rule it hands our LLM-guided outer
loop.

\begin{table}[h]
\centering
\small
\begin{tabular}{@{}p{0.44\linewidth}p{0.5\linewidth}@{}}
\toprule
Cluster resource allocation \citep{cvxcluster2026} & CHEETAH co-design \\
\midrule
Server (a placement slot) & Operation instance (a node at a position in the graph) \\
Shape (servers with identical config) & Unique op-shape (ops with identical dims: identical Pareto/roofline data) \\
Task multi-resource demand (CPU, mem, GPU) & Op demand on the three resources: compute (MACs), traffic (bytes/BW), on-chip capacity (CGM bytes) \\
Bin / server capacity & CGM capacity per timestep, plus the compute roof and the bandwidth roof \\
Binary assignment $x\in\{0,1\}$ & Discrete decisions: tile-config $y$, residency $p$, precision $q$, expert-pin $g$, spec-width $K$ \\
Convex relaxation $x\in[0,1]$ & The LP/SOCP relaxation of the one-hots \\
Anti-affinity as auxiliary capacity rows & Residency, capacity, and McCormick auxiliary constraints \\
Dual variables (dual per resource-per-shape) & Lagrangian duals $\lambda_{\mathrm{CGM}},\lambda_{\mathrm{BW}},\lambda_{\mathrm{MAC}},\lambda_{\mathrm{HBM}},\lambda_{\mathrm{cache}},\lambda_{\mathrm{LD}}$ \\
Net utility $=$ priority $-\sum_r \lambda_r\,\mathrm{demand}_r$ & Reduced cost of a candidate column, $\bar c_j = c_j - \sum_r \lambda_r\, a_{rj}$ \\
Greedy placement by net utility & Dual-guided rounding of the fractional decisions \\
Dynamic re-evaluating as load shifts & Re-solve for updated duals as columns / operating point change \\
Shape aggregation (scale with \#shapes) & Shape aggregation (\cref{app:shape_agg}; $903\to13$ shapes) \\
\bottomrule
\end{tabular}
\caption{One-to-one correspondence between convex cluster resource allocation and
the CHEETAH co-design program.}
\label{tab:cvx_mapping}
\end{table}

\paragraph{Dual variables as the outer-loop search signal.}
The mapping's payoff is that it makes the LLM-guided outer loop
(\cref{sec:llm_outer}) a \emph{dual-guided column generator}, exactly as
\citet{cvxcluster2026} scores task placement. The loop is: (1) solve the inner
SOCP and read the dual dual vector $\lambda=(\lambda_{\mathrm{CGM}},
\lambda_{\mathrm{BW}},\lambda_{\mathrm{MAC}},\dots)$ together with the reduced costs
of the un-chosen columns; (2) the duals \emph{diagnose the bottleneck}, the
largest $\lambda_r$ marks the scarcest resource, while $\lambda_r\approx0$ marks a
free one; (3) the LLM proposes candidate columns (a chip parameterisation, a decision
setting, a tiling) and ranks them by the reduced cost $\bar c_j = c_j -
\sum_r\lambda_r a_{rj}$, i.e.\ objective gain net of the dual-weighted cost of the
resources the column consumes; a column with $\bar c_j<0$ improves the design and
the most negative is the best proposal; (4) re-solve for updated duals and repeat.
The loop terminates when no proposed column has negative reduced cost, a certified
no-improving-candidate stop, with the LP bound at the caps quantifying the maximum
regret of stopping. The LLM only ever \emph{proposes}; the LP scores and certifies,
so no optimality certificate is ever at risk.

\paragraph{Worked example: the duals point to speculative decoding.}
On the DeepSeek-V3 decode workload the converged duals are $\lambda_{\mathrm{BW}}
\gg 0$ (bandwidth is the binding scarcity) and $\lambda_{\mathrm{MAC}}\approx 0$
(compute is idle, hence free). The reduced-cost rule then \emph{selects} the
speculative/parallel-decode-width column $K$ without being told to: verifying $K$
tokens per weight-load spends the near-zero-dual MAC resource to relieve the
high-scored bandwidth (it amortises the mixture-of-experts weight reload over the
accepted tokens), so its reduced cost $\bar c_K = c_K - \lambda_{\mathrm{MAC}}\cdot(\text{K's
compute}) - \lambda_{\mathrm{BW}}\cdot(\text{negative traffic})$ is strongly
negative precisely because $\lambda_{\mathrm{MAC}}\approx0$ while
$\lambda_{\mathrm{BW}}$ is large. The weight-precision column $W4$ is the same logic
on the byte axis (it lowers the $\lambda_{\mathrm{BW}}$-scored traffic directly).
That the dual variables independently point at the two decisions we found empirically
to dominate is evidence the signal is the right one.

\paragraph{Dual-guided rounding.}
The second stage of \citet{cvxcluster2026} also improves our rounding. Where the LP
relaxation leaves a one-hot selector fractional (in our DeepSeek solves the
precision and expert-pin selectors), rounding by reduced cost, fixing first the
selector whose reduced cost most improves the certified bound, then re-solving , 
is tighter than the $\arg\max$ rounding of \cref{sec:v1}, and is the natural
integer companion to the dual-guided search above.

\paragraph{Caveat.}
As in \cref{app:shape_agg}, each round of cluster placement is a static snapshot whereas CHEETAH
carries a time-indexed residency axis; the correspondence is exact on the
shape-invariant block (tiling, roofline, hardware, and the decisions), and the
time-indexed residency is kept per-instance rather than aggregated or scored as a
snapshot.

\section{What CHEETAH-V1 Decides Inside the Program}
\label{app:what_v1_decides}

It is worth stating plainly what is a \emph{decision} of the program and what is
an input to it, because the contribution rests on the boundary between the two.
In a sequential pipeline the hardware is a fixed input: the mapper is asked how
fast a given chip runs a given layer. In CHEETAH-V1 the hardware is a variable,
and the only true constants are the measured Timeloop bases $\Tmin_i(c)$,
$\Tmax_i(c)$ and the per-tile byte footprints. Every coefficient that a
sequential flow would bake in is instead a function of decision variables and is
evaluated inside the constraints.

\paragraph{Hardware, decided rather than given.}
The MAC count and DRAM bandwidth enter through their reference-normalized
reciprocals $v$ and $u$, so a tile's compute time $\Tmin_i(c)\,v$ and memory time
$\Tmax_i(c)\,u$ scale with a chip the program is simultaneously choosing. On-chip
global memory is a one-hot over a size grid, and the HBM residency pool, the
prefix-cache size, the load-engine bandwidth, the systolic array aspect ratio and
the clock and voltage operating point are all selectors in the same program.
These are not swept in an outer loop and read back; they are columns of the
constraint matrix.

\paragraph{Schedule, decided jointly with the hardware.}
Tile configuration per operation, which tensors are resident on chip at each
timestep, which are evicted, which are rematerialized, and which operations fuse,
are decided together with the chip rather than after it. This is what allows a
locally worse tiling to be selected when it produces a globally better schedule,
a trade a map-then-fuse pipeline cannot represent.

\paragraph{Serving behavior, decided as constraints.}
Weight, activation and KV precision are one-hot selectors per phase and per
expert; the KV budget and eviction policy, the mixture-of-experts residency
knapsack, prefix-cache reuse, structured sparsity, memory technology, the
prefill--decode interleaving fraction and the speculative width all enter as
constraint families. Each is inert when pinned, so the program contains the
simpler formulations exactly.

\paragraph{Manufacturability, enforced rather than checked.}
Reticle area, TDP, power density, pin count, logic depth, on-chip interconnect
and SRAM macro decomposition are rows of the program. A design that violates any
of them is infeasible rather than merely flagged, which is why every design in
this paper satisfies all six gates by construction.

\paragraph{What remains outside.}
Accuracy is never linearized: precision entries are admitted only from an
accuracy-vetted menu, and the program chooses among them without modeling
quality. Draft acceptance, the decode batch and the generation length are stated
workload parameters that reshape coefficients but are not decided. Genuinely
discrete or non-convex axes, such as interconnect radix or the draft model
itself, remain columns the outer loop proposes and the program admits only if the
reduced cost is negative and the manufacturability rows clear.

\section{DeepSeek-V3 End-to-End Serving Result}
\label{app:deepseek_summary}

\Cref{tab:deepseek_summary} collects the DeepSeek-V3 mixture-of-experts result in
one place: the end-to-end serving speedup with hardware search inside the loop,
where the time goes, and which mechanism produces the gain. Every symbol is
defined in the table itself so that it can be read without reference to the body.

\paragraph{The serving model.}
End-to-end serving latency is modelled as three phases, not two:
\[
T_{\mathrm{serve}} \;=\; \underbrace{T_{\mathrm{prefill}}}_{\text{context, once}}
\;+\; \frac{G}{A(K)}\Big(
\underbrace{T_{\mathrm{draft}}(K)}_{\text{propose } K \text{ tokens}}
\;+\;
\underbrace{T_{\mathrm{pass}}(K)}_{\text{verify them}}\Big),
\qquad A(K) = \frac{1-p^{K+1}}{1-p},
\]
where $G$ is the number of tokens generated and $A(K)$ is the expected number
accepted per round, so $G/A(K)$ is the number of draft-then-verify rounds. The
draft term is charged rather than assumed free: for an autoregressive draft it
grows as $O(K)$, and for the parallel draft used here it is $O(1)$. At this
operating point the draft is $0.54\%$ of serving, which is small but not zero,
and it is what keeps the selected $K$ finite.

\paragraph{Operating point.}
Three numbers fix the workload. The \emph{decode batch} $B = 32$ is the number of
sequences generated concurrently; it matters because top-$k$ routing makes a
larger batch touch more experts, so it changes the traffic rather than merely
amortizing it. The \emph{draft acceptance} $p = 0.75$ is the probability that a
single token proposed by the draft model is accepted by the target model's
verification; it is a property of the draft/target pair, so we state and sweep it
rather than measure it. The \emph{$1$\,GB diffusion draft} is the small
companion model that proposes candidate tokens: $1$\,GB of weights, and
``diffusion'' meaning it proposes all $K$ tokens in one parallel pass rather than
$K$ sequential ones, so its cost does not grow with the speculative width.

\paragraph{Two questions the selected chip raises.}
The design takes almost no on-chip memory and almost all the bandwidth the
package permits, and both follow from the same fact. In decode the traffic is
$388$\,GB per step of mixture-of-experts weights, which is four orders of
magnitude larger than any on-chip buffer could hold, so additional SRAM has
nothing useful to retain: the weights cannot be cached and the MLA-compressed KV
cache is only $4.3$\,GB. On-chip memory therefore collapses to the smallest grid
point that keeps the schedule feasible, and the area saved is spent elsewhere.
Bandwidth is the opposite case. It divides the dominant term directly, so the
program buys as much as the package allows and stops at $3200$\,GB/s because the
next grid point would need $5250$ signal pins against a $4096$-pin limit. The
resulting design sits at $4000$ pins, which is $97.7\%$ of that limit. This is
the clearest single illustration of why the decisions that matter here reduce or
amortize traffic rather than add storage.

\paragraph{What is solved, and by what.}
The formulation is written as a second-order cone program because two hardware
reciprocals are exact rotated cones. In practice we solve \emph{linear programs}:
the bandwidth reciprocal is already an exact one-hot over a grid, and the
remaining compute reciprocal is eliminated by fixing $v$ on a fine grid, which
makes each grid point a pure LP. Those LPs are solved with HiGHS, which returns a
certified vertex and lets us check every reported design against the constraint
set independently. The Variable-Metric-One-Armed-Rescue solver is used where the
LP is too large or too ill-conditioned for a simplex or interior-point method,
namely the full-resolution instances with per-operation residency, and it agrees
with HiGHS to $10^{-14}$ on the instances where both can run.

\begin{table}[h]
\centering
\footnotesize
\caption{\textbf{DeepSeek-V3 serving co-design with hardware search inside the
loop.}
\emph{Operating point.} Decode batch $B{=}32$ concurrent sequences; draft
acceptance $p{=}0.75$, the probability one proposed token survives verification;
a $1$\,GB parallel (diffusion) draft, which proposes all $K$ tokens in one pass
so its cost does not grow with $K$; context $4096$; $512$ generated tokens.
Speedups are against a fixed TPU-v3 baseline at the same batch. Every row is an
achievable design: a concrete configuration, re-solved and checked against the
constraint set, satisfying all six manufacturability rows.
\emph{Units.} These vary by block and are stated in each block header; speedups
are multiples, so $44.85\times$ means $44.85$ times faster, not $44.85\%$.
\emph{The two speedups answer different questions.} The $44.85\times$ figure is
the best design reachable \emph{given the option menus the program was handed}
(a fixed precision set, a fixed set of memory technologies, $K \in
\{2,8,16,32\}$, and one draft model). The $73.00\times$ figure is what the outer
loop reaches when allowed to \emph{add entries to those menus}, since an
optimizer can only select among options it has been given. An entry is any
option the program could have chosen had it been offered: hardware entries
include a memory technology with a different bandwidth ceiling, area per GB/s
and energy per bit, an omitted bandwidth grid point, another on-chip memory
size, or a different array aspect ratio; software entries include a new
precision format, a sparsity pattern, a speculative width outside the menu, or a
different draft model, which is a paired entry carrying both its size and the
acceptance it achieves. Two were admitted here, one of each kind: a $4$\,GB
adapted draft assumed to reach $p{=}0.90$, and HBM4-class memory. The program
admits an entry only if its reduced cost is negative and the manufacturability
rows still clear.
\emph{Caveats on the larger number.} It is conditional on the assumed $p{=}0.90$,
which we state rather than measure, and on HBM4-class memory being available.
Composition matters: the memory technology is \emph{worse} alone
($49.28\times$) than the draft entry and pays only in combination, which is why
the loop re-solves after each acceptance rather than ranking proposals once.}
\label{tab:deepseek_summary}
\begin{tabular}{@{}llr@{}}
\toprule
\textbf{Quantity} & \textbf{Meaning} & \textbf{Value (unit varies by block)} \\
\midrule
\multicolumn{3}{@{}l}{\textbf{End-to-end serving speedup} \emph{(values are $\times$ faster than TPU-v3, not percentages)}} \\
Best over $K$, fixed menus & best design over the option set the program was & $\mathbf{44.85\times}$ \\
  & given; $K$ is the speculative width (tokens proposed & \\
  & per verification pass), chosen from $\{2,8,16,32\}$ & \\
$+$ outer loop, new options & the loop adds two entries the menus lacked: a & $\mathbf{73.00\times}$ \\
  & $4$\,GB adapted draft assumed to reach $p{=}0.90$ & \\
  & (software) and HBM4-class memory (hardware) & \\
Relaxed bound & no design in this model can beat it & $57.26\times$ \\
\midrule
\multicolumn{3}{@{}l}{\textbf{Where the serving time goes} \emph{(values are \% of total serving latency)}} \\
$T_{\mathrm{prefill}}$ & one compute-bound pass over the $4096$-token context & $42.10\%$ \\
$(G/A(K))\,T_{\mathrm{draft}}$ & $129$ draft rounds proposing $K$ tokens each & $0.54\%$ \\
$(G/A(K))\,T_{\mathrm{pass}}$ & $129$ memory-bound verification passes & $57.35\%$ \\
$A(K)$ & expected tokens accepted per pass at $K^\star{=}16$, $p{=}0.75$ & $3.97$ \\
$G/A(K)$ & number of draft-then-verify rounds for $G = 512$ tokens & $129$ \\
\midrule
\multicolumn{3}{@{}l}{\textbf{What produces the speedup} \emph{(values are cumulative $\times$ speedup, each an achievable design)}} \\
Hardware sizing & $\mathrm{MACs}$, $\mathrm{BW}$, on-chip memory as variables & $4.73\times$ \\
$+$ byte reduction & precision $q_\rho$, $2{:}4$ sparsity, memory technology & $16.49\times$ \\
$+$ speculation & speculative width $K^\star = 16$ & $\mathbf{44.85\times}$ \\
\midrule
\multicolumn{3}{@{}l}{\textbf{Why: where the decode bytes go} \emph{(values are \% of decode DRAM traffic)}} \\
$\Delta^{\mathrm{moe}}$ & mixture-of-experts weight reload, $388.5$ GB/step & $98.6\%$ \\
$\Delta^{\mathrm{KV}}$ & MLA-compressed KV read, $4.29$ GB/step & $1.1\%$ \\
$\Delta^{\mathrm{act}}$ & activations, $1.15$ GB/step & $0.3\%$ \\
\midrule
\multicolumn{3}{@{}l}{\textbf{The chip the program selected} \emph{(values are physical quantities)}} \\
$\mathrm{MACs}$ & INT8 multiply--accumulate units ($24.4\times$ TPU-v3) & $1{,}598{,}953$ \\
$\mathrm{BW}$ & DRAM bandwidth & $3200$ GB/s \\
$C_{\mathrm{GM}}$ & on-chip global memory & $1.0$ MiB \\
$n_{\mathrm{PE}}$ & processing elements, $= \lceil \mathrm{MACs}/(s_x s_y)\rceil$ & $98$ \\
$s_x \times s_y$ & systolic array dimensions per PE (rows $\times$ columns), & $128 \times 128$ \\
  & giving $16{,}384$ MACs per processing element & \\
DRAM channels & at ${\sim}100$ GB/s per channel & $32$ \\
\midrule
\multicolumn{3}{@{}l}{\textbf{Manufacturability: six of six rows satisfied} \emph{(achieved value vs.\ cap)}} \\
Pin count & $\mathrm{BW}$ at $0.8$ GB/s per HBM3 pin (cap $4096$) & $4000$ \\
Reticle area & die area (cap $600$ $\mathrm{mm}^2$) & $497.3$ $\mathrm{mm}^2$ \\
Power density & dynamic $320$ W over the die (cap $2.0$) & $0.64$ W/$\mathrm{mm}^2$ \\
Logic depth & gate levels at $1$ GHz (cap $35$) & $25$ \\
Crossbar & PE-to-bank routing cells (cap $5{\times}10^6$) & $1568$ \\
SRAM macros & $64$ KiB banks (cap $4096$) & $16$ \\
\bottomrule
\end{tabular}
\end{table}

The decomposition explains itself. Decode dominates serving, and decode is
almost entirely mixture-of-experts weight traffic, so the mechanisms that reduce
or amortize that traffic are the ones that matter: sizing bandwidth and on-chip
memory first, then reducing the bytes through precision and sparsity, then
amortizing each remaining weight load over more accepted tokens through
speculation. Tiling and fusion contribute by keeping the prefill phase on its
compute floor, but they cannot move the decode wall, which is why a formulation
that decides tiling alone plateaus and one that also sizes the hardware and
decides the serving decisions does not.

\section{A Worked Example: Hardware Coefficients as Functions, Not Constants}
\label{app:worked_example}

The central modelling decision in this paper is that hardware quantities enter
the program as functions of decision variables rather than as fixed
coefficients. It is worth seeing exactly what that changes on a single concrete
quantity.

\paragraph{The quantity.}
Every decode step must read the routed mixture-of-experts weights from DRAM.
Timeloop measures $388.5$\,GB per step at $8$-bit precision. That byte count is a
genuine constant: it is a property of the model and its routing, not of any chip.
The \emph{time} it takes is bytes divided by bandwidth, and that is where the two
formulations part company.

\paragraph{As a constant, in a sequential pipeline.}
Fix the chip first at the TPU-v3 bandwidth of $900$\,GB/s. The coefficient is now
a number,
\[
t_{\mathrm{decode}} \;=\; \frac{388.5\ \mathrm{GB}}{900\ \mathrm{GB/s}}
\;=\; 0.432\ \mathrm{s},
\]
and $0.432$ enters the objective as a literal. Asking what a different bandwidth
would buy requires discarding the program and building another one.

\paragraph{As a function, in CHEETAH-V1.}
The corresponding row is
\[
\tau \;\ge\; \underbrace{\Tmax(c)}_{\text{measured bytes}}
\;\times\; \underbrace{s^W}_{\text{precision variable}}
\;\times\; \underbrace{u}_{\text{bandwidth variable}},
\]
where $s^W$ is set by the precision one-hot ($1.0$ at W8, $0.5$ at W4) and
$u = \mathrm{BW_{ref}}/\mathrm{BW}$. Neither factor is a number until the solver
chooses one.

\paragraph{Why this matters: the two decisions multiply.}
\Cref{tab:worked_example} evaluates the same row at the four combinations of
these two decisions. The best cell is $7.11\times$ faster than the baseline, and
that factor decomposes as $2\times$ from precision and $3.56\times$ from
bandwidth: the decisions compose rather than compete.

\begin{table}[h]
\centering
\footnotesize
\caption{Decode weight-streaming time for one step of DeepSeek-V3, evaluated at
each combination of the precision and bandwidth decisions. A pipeline that fixes
bandwidth before considering precision is confined to the left column.}
\label{tab:worked_example}
\begin{tabular}{@{}lcc@{}}
\toprule
 & $\mathrm{BW} = 900$\,GB/s & $\mathrm{BW} = 3200$\,GB/s \\
\midrule
W8 \quad($388.5$\,GB) & $0.432$\,s \quad (baseline) & $0.121$\,s \\
W4 \quad($194.3$\,GB) & $0.216$\,s & $\mathbf{0.061}$\,\textbf{s} \\
\bottomrule
\end{tabular}
\end{table}

The consequence is stronger than a search-quality argument. A pipeline that fixes
bandwidth at $900$\,GB/s before precision is considered can only ever reach the
left column, so its best attainable answer is $0.216$\,s. The $0.061$\,s cell is
not merely missed by a weak search; it is \emph{unreachable}, because by the time
precision is being chosen the bandwidth has already become a constant.

\paragraph{What it costs.}
Reading the row once more, $s^W \times u$ is a product of two decision variables.
It is bilinear and therefore nonconvex, which is exactly why the formulation
carries a McCormick envelope at that position, and why $u$ requires the rotated
cone $\mathrm{BW_{eff}}\,u \ge \mathrm{BW_{ref}}$. The multiplicative gain in
\cref{tab:worked_example} and the solver machinery of \cref{sec:solver_vm} are
the same fact viewed from opposite sides: expressiveness bought with
nonconvexity, and nonconvexity paid for with envelopes, cones, and a solver that
tolerates the resulting conditioning.

This is not a hypothetical. It is the structure that produced the reported
design, which buys $3200$\,GB/s, the most the pin-count row permits, and selects
W4 with a $4$-bit KV cache: both halves of the bottom-right cell.

% ---------------------------------------------------------------------------
\section{Where Each Workload Family Slows Down, and the Local Sacrifice That Wins}
\label{app:family_bottlenecks}

The formulation earns its keep only where a \emph{locally} worse decision is
\emph{globally} better, because that is precisely the trade a sequential
pipeline cannot represent: a mapper that has already committed to the
per-operation optimum has priced the local choice and thrown away the coupling
that would have paid for it. \Cref{tab:family_bottlenecks} names, for each
workload family in our suite, where inference actually slows down and which
locally costly choice the joint program is able to buy in exchange. Read it as
the qualitative statement of which term in the objective the coupling is
expected to move; the measured per-family decompositions are in
\cref{app:deficit,app:blueprints}.

\begin{table}[ht!]
\centering
\small
\caption{Per-family inference bottlenecks and the locally costly choice that can
win globally. Each row is an instance of the same structural claim: the
per-operation optimum and the end-to-end optimum disagree, and the size of the
disagreement is what the joint program recovers. The EfficientNet-B7 figure is
FAST's own profile on TPU-v3~\citep{zhang_fast}, where depthwise operations
account for $5\%$ of FLOPs but $65.3\%$ of runtime, which is the clearest
published case of a bottleneck that FLOP-weighted reasoning cannot see.}
\label{tab:family_bottlenecks}
\begin{tabularx}{\linewidth}{@{}p{0.17\linewidth}XX@{}}
\toprule
\textbf{Family} & \textbf{Where inference slows down} &
\textbf{Locally costly choice that can win globally} \\
\midrule
Llama-1 13B/65B, Llama-3 70B
& Prefill stresses projection/FFN GEMMs and attention. Small-batch decode often
stresses weight bandwidth; longer contexts increase KV traffic and attention
costs.
& A smaller attention or FFN tile preserves weights or KV for subsequent
operations. Splitting attention across the KV history adds reduction work but
can improve parallelism. \\
\addlinespace
EfficientNet B1--B7
& Depthwise utilization, expanded activations, SE reductions, and layout/halo
traffic. FAST's historical B7 profile found depthwise operations were only $5\%$
of FLOPs but $65.3\%$ of runtime on TPU-v3.
& A slower compact tile preserves an expanded feature or residual, avoiding a
spill and reload. \\
\addlinespace
SD1.5 UNet
& Convolutions, attention, and encoder skips that remain live until much later
decoder operations. Complete image generation repeats the denoiser.
& Smaller intervening tiles preserve a valuable skip; legal recomputation can
release memory for something more expensive to reload. \\
\addlinespace
DeepSeek-V3 MoE
& Expert-weight movement, small or uneven expert GEMMs, MLA attention, and, when
distributed, dispatch/combine communication.
& Compact mappings preserve experts or KV; prefetching and grouping incur costs
now to avoid larger later stalls. \\
\bottomrule
\end{tabularx}
\end{table}

Two properties of the table are worth making explicit. First, every right-hand
entry is a decision the V1 program can express as a column and a sequential
pipeline cannot, because in each case the locally costly choice is only
justified by a constraint that binds at a \emph{different} operation or a
\emph{later} timestep: the tile that looks slow in isolation is paying for a
residency that a downstream operation will consume. Second, the families differ
in \emph{which} resource the sacrifice buys, on-chip capacity for EfficientNet
and SD1.5, memory bandwidth for the Llamas and DeepSeek-V3, which is why the
same formulation produces qualitatively different chips across the suite rather
than one design rescaled.

% ---------------------------------------------------------------------------
\section{A Minimal Certified Example of Hardware/Software Non-Separability}
\label{app:synergy_example}

\Cref{tab:family_bottlenecks} asserts that a locally costly choice can win
globally. This appendix gives the smallest instance we know of in which that
claim is not only true but \emph{certified}, and in which the failure mode it
warns against, alternating hardware-only and software-only optimization, can be
exhibited exactly. Every quantity is a synthetic normalized unit; this is a
teaching model with a closed form, not a chip, a workload, or a measurement.

\paragraph{Setting.}
Two serial invocations of one kernel share an immutable weight. Two functionally
equivalent implementations are admitted: a \emph{fast} kernel (compute $1.0$,
scratch $40$) and a \emph{compact} kernel (compute $1.2$, scratch $20$). The
weight occupies $20$ SRAM units if retained between invocations; its first-use
load costs $10$ traffic units and a second-use miss costs $10$ more. A chip of
total area $100$ is divided continuously among compute $m$, bandwidth $b$, and
SRAM $S$, with $S$ at least $40$ so that both kernels fit every admitted chip.
For a branch with compute demand $C$, traffic $Q$, and required SRAM $S$, the
latency $C/m + Q/b$ subject to $m+b=100-S$ has the closed-form optimum
%
\begin{equation}
\label{eq:synergy_closed_form}
T^\star=\frac{(\sqrt{C}+\sqrt{Q})^{2}}{100-S},\qquad
m^\star=(100-S)\,\frac{\sqrt{C}}{\sqrt{C}+\sqrt{Q}} .
\end{equation}

\paragraph{The four moves.}
Eight branches (two kernels per invocation, retain or not) are all covered by
\cref{eq:synergy_closed_form}. \Cref{tab:synergy_moves} records the four moves an
optimizer can make from the locally natural starting point.

\begin{table}[ht!]
\centering
\small
\caption{Alternating optimization stalls; the joint move does not. Latencies in
normalized units from \cref{eq:synergy_closed_form}. The fast kernel is the
per-kernel latency optimum on any chip with adequate SRAM, which is what a
per-operation mapper selects; ``locally fastest'' is that mapper's choice, not
a claim that it is globally wise.}
\label{tab:synergy_moves}
\begin{tabular}{@{}lccc@{}}
\toprule
\textbf{Move} & \textbf{Kernel} & \textbf{Chip} & \textbf{Latency} \\
\midrule
Chip optimized for the locally fastest kernel & fast & optimal for fast & $0.5236$ \\
Change software only, on that chip & compact & unchanged & $0.5560$ \\
Re-optimize hardware only, keeping the fast kernel & fast & re-optimized & $0.5236$ \\
\textbf{Change hardware and software jointly} & compact & re-optimized & $\mathbf{0.3700}$ \\
\bottomrule
\end{tabular}
\end{table}

Both single-axis moves are stuck: the fast-kernel chip is already optimal for
the fast kernel, and on that chip the compact kernel is strictly worse. Only the
joint move improves, by $1.415\times$, and the mechanism is exactly the one
\cref{tab:family_bottlenecks} describes for EfficientNet and SD1.5: the compact
kernel needs $20$ fewer SRAM units, which lets the weight be retained
\emph{within} the $40$-unit floor, and the released area is re-spent on compute
and bandwidth ($m$ from $12.36$ to $19.73$, $b$ from $27.64$ to $40.27$). The
slower kernel is worth choosing only after the chip has been resized to exploit
it, and the chip is worth resizing only if the slower kernel will be chosen.
Neither decision is justified by the other's status quo, which is why a
sequential pipeline cannot reach the point.

\paragraph{The result is robust, not a knife edge.}
The $40$-unit SRAM floor coincides with the fast kernel's scratch, so the
compact kernel's retention is exactly free at that value; one might suspect the
effect is tuned. It is not. Sweeping the floor over its admissible range, the
joint move wins throughout:
%
\begin{center}
\small
\begin{tabular}{@{}lcccc@{}}
\toprule
SRAM floor & $40$ (shipped) & $45$ & $50$ & $30$--$35$ \\
\midrule
Best fast-kernel design & $0.5236$ & $0.5236$ & $0.5236$ & $0.5236$ \\
Joint (compact $+$ retain) & $0.3700$ & $0.4036$ & $0.4440$ & $0.3700$ \\
Ratio & $1.415$ & $1.297$ & $1.179$ & $1.415$ \\
\bottomrule
\end{tabular}
\end{center}
%
Below $40$ the fast kernel no longer fits in isolation, so the comparator
changes meaning; above it the advantage shrinks smoothly and never vanishes.

\paragraph{The bound is a certificate, not a restatement.}
Each branch is also solved as an LP in which $C/m$ and $Q/b$ are replaced by
their supporting planes at $(m^\star,b^\star)$, and HiGHS returns exactly
$T^\star$ for every branch. Because a tangent plane of a convex function is a
global underestimator, that LP value is a valid lower bound on the branch
optimum for \emph{any} linearization point; it equals the replayed latency at
the linearization point only when that point is optimal. We verified this has
teeth by perturbing the linearization point: at $\pm1\%$ the LP bound falls to
$0.3664$ against a replay of $0.3700$, at $\pm10\%$ to $0.336$ against $0.372$,
and equality is recovered only at $m^\star$. An LP whose optimum is attained at
its own linearization point is the first-order optimality certificate for a
convex program, so the eight branch optima are certified rather than assumed,
and the $1.415\times$ is a gap between two certified values.

\paragraph{What it does and does not show.}
It shows that the objective is non-separable in hardware and software even in
the simplest case with continuous ideal resources, and that this
non-separability defeats coordinate-wise optimization. It does not model
tiling, time-indexed residency, array parallelism limits, or any physical
constraint, and its $1.415\times$ carries no information about DeepSeek-V3 or
any chip; the measured per-family decompositions are in
\cref{app:deficit,app:blueprints}. Its value is that the certified argument is
short enough to check by hand.
